\newpage
\section{Lesson 5, 01/10}
\subsection{More on the Lorentz group}
The generic Lorentz transformation can also be written in terms of $J,K$:
$$\Lambda = \exp\left[i\vartheta_iJ^i+i\eta_iK^i\right]$$
where $\vartheta_i$ are the Euler angles (parameters of the rotation) and $\eta_i$
are the rapidities (parameters of the boost).

With some work, it is possible to prove that
$$\frac{1}{2}M_{\mu\nu}M^{\mu\nu} = J_iJ^i - K_iK^i = \mathbf{J^2}-\mathbf{K^2}$$
and
$$-\frac{1}{4} \varepsilon^{\mu\nu\rho\sigma}M_{\mu\nu}M_{\rho\sigma} = J_iK^i = 
\mathbf{J}\cdot\mathbf{K}$$
are Casimir operators for the representation: there exist only 2 linear independent
Casimir operators and each and any linear combination of these two is valid, but
these two are important because they are, respectively, symmetric and antisymmetric
(how can they be?? They're scalar!).

\textbf{Theorem:} The algebra of the Lorentz Group can be factorized as the direct
product of two $SU(2)$ algebras
$$\mathfrak{so}(1,3) = \mathfrak{su}(2)\otimes\mathfrak{su}(2)$$

This suggests to find new generators that are suited to this decomposition
\begin{gather*}
\mathbf{N} = \frac{1}{2}\left(\mathbf{J}+i\mathbf{K}\right)\\
\mathbf{N}^\dagger = \frac{1}{2}\left(\mathbf{J}-i\mathbf{K}\right)
\end{gather*}
their algebra is
\begin{align*}
&\left[N_i,N_j^\dagger\right] = 0\\
&\left[N_i,N_j\right] = i\varepsilon_{ijk}N^k\\
&\left[N_i^\dagger,N_j^\dagger\right] = i\varepsilon_{ijk}N^{k\dagger}
\end{align*}
\textbf{Observation:} These generators are convenient, since the algebra
of the $N$'s and that of the $N^\dagger$'s are closed and disjoint.

The Casimir operators are
\begin{align*}
    N_iN^i &\longrightarrow n(n+1)\\
	   &\text{\footnotesize eigenvalues}\\
    N^\dagger_iN^{i\dagger} &\longrightarrow m(m+1)    ,
\end{align*}
so the irreducible representations of the group will be labeled by the pair $(n,m)$.
Since they obey $SU(2)$ algebra, it is clear that $n,m$ will have integer or
semi-integer values and that the eigenvalues of
$N_3,N_3^{\dagger}$, $s$ and $s^\dagger$, will respect
$$-n<s<n\qquad-m<s^\dagger<m.$$
To each pair $n,m$ corresponds a name for the field
\begin{table}[ht!]
    \centering
    \begin{tabular}{c|c|c|c}
	$(n,m)$&Field name&Dim&Example\\ \hline
	$(0,0)$&Scalar $\varphi$&1&Higgs \\
	$(1/2,0)$&Spinor $\psi_L$&2&Weyl Spinor\\
	$(0,1/2)$&Spinor $\psi_R$&2&Weyl Spinor\\
	$(1/2,0)\oplus(0,1/2)$&Dirac Spinor&4&\\
	$(1/2,1/2)$&4-vector field (spin 1)&4&$A^\mu$ ELM\\
	$(1,0)\oplus(0,1)$&Rank 2 antisymm tensor&6&$F^{\mu\nu}$ ELM
    \end{tabular}
\end{table}
\newpage
\subsection{Poincaré Group}
The QFT must also be invariant over spacetime \textit{translations}: the Poincaré
group includes Lorentz transformations \textit{and} translations
$$x^\mu\rightarrow x'^\mu:=\Lambda\indices{^\mu_\nu}x^\nu + a^\mu$$
the group parameters are $6+4=10$. The generators of the 4-translations,
similarly to QM, are the 4-momenta
$$P^\mu = (H,P^i);$$
moreover, in QM1 the time evolution operator was found to be
$$U=e^{-iH}$$
which, in hindsight, is an exponential map, in which $H$ is the generator and $t$ is
the parameter!

There will be some problems with the fact that the translation group is not
compact.

Finding the commutation relations for the algebra is not trivial, since rotations
and translations do not commute, even in Minkowski space. Consider two
successive Poincaré transformations (why????):
$$x^\mu\rightarrow{(\Lambda_{(1)})}\indices{^\mu_\nu}x^\nu + a_1^\mu \rightarrow
{(\Lambda_{(2)})}\indices{^\mu_\rho}{(\Lambda_{(1)})}\indices{^\mu_\nu}x^\nu + 
{(\Lambda_{(2)})}\indices{^\mu_\rho}a_1^\rho+a_2^\mu$$
The commutation relations are
\begin{align*}
&\left[P^\mu,P^\nu\right]=0\\
&\left[M_{\mu\nu},P_\rho\right] = -ig_{\rho\sigma}P_\nu+ig_{\nu\rho}P_\mu\\
\end{align*}
They are clearer when time and space indexes are expressed separately
\begin{align*}
&\left[P^i,P^j\right]=0\\
&\left[P^i,H\right]=0\\
&\left[K^i,P^j\right]=i\delta^{ij}H\\
&\left[J^i,J^j\right]=i\varepsilon\indices{^i^j_k}J^k\\
&\left[K^i,H\right]=iP^i\\
&\left[J^i,P^j\right]=i\varepsilon\indices{^i^j_k}P^k\\
&\left[J^i,H\right]=0\\
&\left[K^i,K^j\right]=i\varepsilon\indices{^i^j_k}J^k\\
&\left[J^i,K^i\right]=i\varepsilon\indices{^i^j_k}K^k\\
\text{REORDER}
\end{align*}
Next, the Casimir operators: a natural choice is
$$P^2 = P_\mu P^\mu$$
a not so natural choice is the \textit{Pauli-Lubanski vector} squared:
$$W^\mu = \varepsilon^{\mu\nu\rho\sigma}P_\nu M_{\rho\sigma}\rightarrow W^2 = W_\mu
W^\mu.$$
A nice relation is
$$P_\mu W^\mu = 0$$
for antisymmetry reasons.

